\documentclass[10pt,a4paper]{amsart}
\usepackage[margin=19mm]{geometry}
\usepackage{amsmath,amssymb,mathtools}
\usepackage[scr=rsfs,cal=euler]{mathalfa}
\usepackage{xfrac}
\usepackage[unicode,hidelinks,bookmarks=false]{hyperref}
\newcommand{\ie}{i.e.\ }
\newcommand{\cf}{cf.\ }
\newcommand{\resp}{respectively\ }
\newcommand{\Subsection}{\subsection}
\providecommand{\nobreakdash}{\nobreak\hbox{-}}
\setlength{\emergencystretch}{2em}
\multlinegap0pt
\usepackage{bm}
\usepackage{amssymb}
\usepackage[shortlabels]{enumitem}
\usepackage{xcolor}
\usepackage[normalem]{ulem}
\newtheorem{theorem}{Theorem}
\newtheorem{proposition}{Proposition}
\usepackage{cleveref}
\crefname{theorem}{Theorem}{Theorems}
\crefname{proposition}{Proposition}{Propositions}
\renewcommand{\a}{\alpha}
\renewcommand{\b}{\beta}
\renewcommand{\l}{\lambda}
\newcommand{\os}[2]{\overset{#1}{#2}}
\renewcommand{\pmod}[1]{\ \left( \mathrm{mod} \, #1 \right)}
\renewcommand{\r}{{\rho}}
\renewcommand{\t}{\tau}
\newcommand{\z}{\zeta}
\title{Proof of a conjecture of Andrews and Bachraoui on a Hecke sum}
\author{Koustav Banerjee, Kathrin Bringmann}
\date{Source: arXiv:2605.10300v1 (CC BY 4.0)}
\begin{document}
\maketitle

\noindent\emph{Source and licence.} Proof of a conjecture of Andrews and Bachraoui on a Hecke sum. Version arXiv:2605.10300v1 (CC BY 4.0), distributed by arXiv under the Creative Commons Attribution 4.0 International licence, http://creativecommons.org/licenses/by/4.0/, as recorded in the arXiv licence metadata for this exact version and shown on the abstract page https://arxiv.org/abs/2605.10300v1. Full licence notice: \texttt{licenses/arxiv-2605.10300-CC-BY-4.0.txt}.\par\smallskip

\noindent\textit{Editorial scope.} Counted unit: the article's proposition hathtrans (source0.tex lines 817--822). The article's complete original proof of it is delivered with it (source0.tex lines 823--897, one \texttt{proof} environment of 75 lines). No mathematics is written, repaired or re-derived: the statement and the proof are the article's own lines, in the article's own notation, and no retained line is substituted or re-flowed.  Retained uncounted as the source-local context the proof needs: lines 415--427; lines 442--447; lines 488--501.  Nothing was omitted from any retained span, so the package carries no omission record.  No retained line contains a citation, so the wrapper carries no bibliography and no bibliography entry.  No retained line contains a figure environment, an image-inclusion command or drawing code, so no image is delivered and the package declares no asset dependency.  Editorial formatting only, in addition: an AMS \texttt{amsart} preamble that restates the article's own declarations and macros used by the retained lines, namely the environments \texttt{theorem}, \texttt{proposition} on separate counters, together with the standard packages those lines need.  The retained environments print in this document as Theorem 1, Proposition 1 in order of appearance, whereas the article prints the same items with the numbering of its own full document.  The complete native source of the article as distributed in the arXiv e-print for this exact version is bundled unchanged as \texttt{sources/200288/source0.tex}.
\begin{theorem}\label{Zwmainthm1}
	The function $\vartheta_{\bm{a},\bm{b}}$ has the following elliptic and modular transformation properties:
	\begin{enumerate}
		\item $\vartheta_{\bm{a}+\bm{\l},\bm{b}}=\vartheta_{\bm{a},\bm{b}}$ for $\bm{\l}\in \mathbb{Z}^{\ell}$,
		\item $\vartheta_{\bm{a},\bm{b}+{\bm \mu}}=e^{2\pi iB(\bm{a},\bm{\mu})}\vartheta_{\bm{a},\bm{b}}$ for $\bm{\mu}\in A^{-1}\mathbb{Z}^{\ell}$,
		\item $\vartheta_{-\bm{a},-\bm{b}}=-\vartheta_{\bm{a},\bm{b}}$,
		\item $\vartheta_{\bm{a},\bm{b}}(\t+1)=e^{-2\pi i Q(\bm{a})-\pi i B(A^{-1}A^*,\bm{a})}\vartheta_{\bm{a},\bm{b}+\bm{a}+\frac 12A^{-1}A^{*}}(\t)$ with $A^*$ denoting the vector of diagonal elements of $A$.
		\item If $\bm{a}, \bm{b}\in R(\bm{c_1})\cap R(\bm{c_2})$, then
		\[
		\vartheta_{\bm{a},\bm{b}}\left(-\frac 1\t\right)=\frac{i}{\sqrt{-\det (A)}}\left(-i\t\right)^{\frac{\ell}{2}}e^{2\pi i B(\bm{a},\bm{b})}\sum_{\bm{\nu}\in A^{-1}\mathbb{Z}^{\ell}\pmod{\mathbb{Z}^{\ell}}}\vartheta_{\bm{b}+\bm{\nu},-\bm{a}}(\t).
		\]
	\end{enumerate}
\end{theorem}

\begin{theorem}\label{Zwmainthm2}
	Let $C\in O^{+}_A(\mathbb{Z})$, $\bm{c_1}, \bm{c_2}\in \overline{C}_Q$, and $\bm{a}\in R(\bm{c_1})\cap R(\bm{c_2})$. Then
	\[
	\vartheta^{C\bm{c_1}, C\bm{c_2}}_{C\bm{a},C\bm{b}}(\t)=\vartheta^{\bm{c_1}, \bm{c_2}}_{\bm{a},\bm{b}}(\t).
	\]
\end{theorem}

\begin{proposition}\label{prop1}
We have
\[
\widehat{H}(\t)=\frac{e^{-\frac{\pi i}{3}}}{2}\vartheta_{\bm{a},\bm{b}}(\t).
\]
%\[
%H(q)=\frac{e^{-\frac{\pi i}{3}}}{2}q^{-\frac 23}\!\!\sum_{\bm{n}\in \mathbb{Z}^2+\bm{a}}\left(\textnormal{sgn}\left(B\left(\bm{c_1},\bm{n}\right)\right)-\textnormal{sgn}\left(B\left(\bm{c_2},\bm{n}\right)\right)\right)e^{2\pi i B\left(\bm{n}, \bm{b}\right)}q^{2Q\left(\bm{n}\right)}
%\]
%and its non-holomorphic part is
%\begin{align}\label{shadowH}
%&-\frac{q^{-\frac 23}}{\sqrt{\pi}}\sum_{\varepsilon\in\{0,1\}}\!\!(-1)^{\varepsilon}\sum_{\substack{n\in \mathbb{Z}\\r\in \mathbb{Z}}}\textnormal{sgn}\left(2r+\varepsilon+\frac 13\right)\nonumber\\
%&\hspace{1.5 cm}\times \Gamma\left(\frac 12,12\pi\left(2r+\varepsilon+\frac 13\right)^2v\right)q^{-3\left(2r+\varepsilon+\frac 13\right)^2+\left(2n+\varepsilon+1\right)^2}.
%\end{align}
\end{proposition}

\begin{proposition}\label{hathtrans}
We have\footnote{Again we could prove that $\widehat{H}$ is a harmonic Maass form of weight $1$ with multiplier.}
\[
\widehat{H}\left(\t+1\right)=\z_3\widehat{H}(\t),\qquad \widehat{H}\left(\frac{\t}{2\t+1}\right)=\z_{12}(2\t+1)\widehat{H}(\t).
\]
\end{proposition}

\begin{proof}
By Theorem \ref{Zwmainthm1} (4), we have 
\[
\vartheta_{\bm{a},\bm{b}}(\t+1)=\z^{2}_3\vartheta_{\bm{a},\bm{b}+\left(\frac{5}{6},\frac 12\right)}(\t).
\]
Using Theorem \ref{Zwmainthm1} (2) and Proposition \ref{prop1} gives the first claim. Note that the shift is allowed since $(\frac 56,\frac 12)$ lies in $A^{-1}\mathbb{Z}^2$.

%Now we use {\bf KB: I think we only use (1).} \Cref{Zwmainthm1} (1) and \Cref{Zwmainthm1} (2). 

We next note that, by Theorem \ref{Zwmainthm1} (5)
\begin{align}\label{star}
\vartheta_{\bm{a},\bm{b}}\left(-\frac{1}{\t}\right)%&=\frac{i}{\sqrt{-\det(A)}}(-i\t)e^{2\pi i B\left(\bm{a},\bm{b}\right)}\sum_{\bm{\nu}\in A^{-1}\mathbb{Z}^2\pmod{\mathbb{Z}^2}}\vartheta_{\bm{b}+\bm{\nu},-\bm{a}}(\t)\nonumber\\
&=\frac{\z_6\t}{2\sqrt{3}}\sum_{\nu_1,\nu_2}\vartheta_{\bm{b}+\bm{\nu},-\bm{a}}(\t),
\end{align}
where $\nu_1$ runs through a set of representatives of $(\frac 16\mathbb{Z})/\mathbb{Z}$ and $\nu_2$ runs through a set of representatives of $(\frac 12\mathbb{Z})/\mathbb{Z}$. We substitute $\nu_1=\frac{\ell_1}{6}$ and $\nu_2=\frac{\ell_2}{2}$ with $\ell_1$ running modulo $6$ and $\ell_2$ running modulo $2$. We then define
\begin{align}\label{l1l2}
\bm{\rho_{\ell_1,\ell_2}}&:=\bm{b}+\left(\frac{\ell_1}{6},\frac{\ell_2}{2}\right).%=\left(\frac{1}{12},-\frac 14\right)+\left(\ell_1,\ell_2\right)\os{\substack{\ell_1\mapsto\frac{\ell_1}{6}\\\ell_2\mapsto \frac{\ell_2}{2}}}=\left(\frac{1}{12}+\frac{\ell_1}{6},-\frac 14+\frac{\ell_2}{2}\right)\\
%&=\left\{\left(\frac{1}{12},-\frac 14\right), \left(\frac{1}{4},-\frac 14\right), \left(\frac{5}{12},-\frac 14\right), \left(\frac{7}{12},-\frac 14\right), \left(\frac{3}{4},-\frac 14\right),\right.\nonumber\\
%&\hspace{0.5 cm}\left. \left(\frac{11}{12},-\frac 14\right), \left(\frac{1}{12},\frac 14\right),\left(\frac{1}{4},\frac 14\right), \left(\frac{5}{12},\frac 14\right), \left(\frac{7}{12},\frac 14\right), \left(\frac{3}{4},\frac 14\right), \left(\frac{11}{12},\frac 14\right)\right\}.\nonumber
\end{align}
Now let $C:=(\begin{smallmatrix}
-1&0\\
\ 0&1
\end{smallmatrix})$. Note that  $C\in O^{+}_{A}(\mathbb{Z})$ as $C^TAC=A$ and $C\bm{c_1}=\bm{c_2}$ which lies in the same component as $\bm{c_1}$. Thus, by Theorem \ref{Zwmainthm2}, we have
\begin{equation*}
\vartheta^{\bm{c_1}, \bm{c_2}}_{\bm{\a},\bm{\b}}=\vartheta^{C\bm{c_1}, C\bm{c_2}}_{C\bm{\a},C\bm{\b}}.
\end{equation*}
Noting that $C\bm{c_1}=\bm{c_2}$, $C\bm{c_2}=\bm{c_1}$, we obtain, as interchanging $\bm{c_1}$ and $\bm{c_2}$ changes the sign of the theta function, 
\[
\vartheta_{\bm{\a},\bm{\b}}=-\vartheta_{C\bm{\a},C\bm{\b}}.
\]
Consequently, we obtain
\begin{equation}\label{plug1}%\label{c1c2}
\vartheta_{\bm{\rho_{\ell_1,\ell_2}},-\bm{a}}=-\vartheta_{C\bm{\rho_{\ell_1,\ell_2}},C(-\bm{a})}=-\vartheta_{C\bm{\rho_{\ell_1,\ell_2}},\bm{a}},
\end{equation}
 using the fact that $C(-\bm{a})=\bm{a}$ in the final step. Now, by \eqref{l1l2}, we have 
\[
C\bm{\rho_{\ell_1,\ell_2}}\equiv \bm{\rho_{5-\ell_1,\ell_2}}\pmod{\mathbb{Z}^2}.
\]
Plugging this into \eqref{plug1} and using Theorem \ref{Zwmainthm1} (2) gives 
\[
\vartheta_{\bm{\rho_{\ell_1,\ell_2}},-\bm{a}}%=-e^{2\pi iB\left(\bm{\rho_{5-\ell_1,\ell_2}},2\bm{a}\right)}\vartheta_{\bm{\rho_{5-\ell_1,\ell_2}},-\bm{a}}
=-\z^{\ell_1+2}_3\vartheta_{\bm{\rho_{5-\ell_1,\ell_2}},-\bm{a}}.
\]
 Inserting this into \eqref{star} yields 
\begin{align}\label{ref3}
\vartheta_{\bm{a},\bm{b}}\left(-\frac{1}{\t}\right)&=\frac{\z_6\t}{2\sqrt{3}}\sum_{\substack{\ell_1\in\{0,1,2\}\\\ell_2\in\{0,1\}}}\left(1-\z^{2\ell_1+1}_3\right)\vartheta_{\bm{\rho_{\ell_1,\ell_2}},-\bm{a}}(\t)\nonumber\\
&=\frac{\z_6\t}{2\sqrt{3}}\sum_{\substack{\ell_1\in\{0,2\}\\\ell_2\in\{0,1\}}}\left(1-\z^{2\ell_1+1}_3\right)\vartheta_{\bm{\rho_{\ell_1,\ell_2}},-\bm{a}}(\t),
\end{align}
as the term with $\ell_1=1$ vanishes.

 Next, by Theorem \ref{Zwmainthm1} (4), we obtain 
\begin{align*}%\label{ref1}
&\vartheta_{\bm{\rho_{\ell_1,\ell_2}},-\bm{a}}(\t-2)\nonumber\\
&\hspace{1.5cm}=e^{4\pi i Q\left(\bm{\rho_{\ell_1,\ell_2}}\right)+2\pi i B\left(A^{-1}A^*,\bm{\rho_{\ell_1,\ell_2}}\right)}\vartheta_{\bm{\rho_{\ell_1,\ell_2}},-\bm{a}-2\bm{\rho_{\ell_1,\ell_2}}-A^{-1}A^*}(\t).
\end{align*}
Using this, \eqref{l1l2}, and Theorem \ref{Zwmainthm1} (2), we have (as the relevant shift is $-2\bm{\rho_{\ell_1,\ell_2}}-A^{-1}A^{*}$ which lies in $A^{-1}\mathbb{Z}^2$)
\begin{equation}\label{ref4}
\vartheta_{\bm{\rho_{\ell_1,\ell_2}},-\bm{a}}(\t-2)=e^{-4\pi i Q\left(\bm{\rho_{\ell_1,\ell_2}}\right)}\vartheta_{\bm{\rho_{\ell_1,\ell_2}},-\bm{a}}(\t).
\end{equation}
Now, by \eqref{l1l2}, we have 
\begin{align*}
Q\left(\bm{\rho_{\ell_1,\ell_2}}\right)%\os{\eqref{l1l2}}=3\left(\frac{1}{12}+\frac{\ell_1}{6}\right)^2-\left(-\frac 14+\frac{\ell_2}{2}\right)^2
=\begin{cases}
-\frac{1}{24}\ \ &\text{if}\ \ \ell_1=0, \ell_2\in\{0,1\},\\
\frac{11}{24}\ \ &\text{if}\ \ \ell_1=2, \ell_2\in\{0,1\}.
\end{cases}
\end{align*}
Plugging this into \eqref{ref4} gives 
\begin{equation*}%\label{ref5}
\vartheta_{\bm{\rho_{\ell_1,\ell_2}},-\bm{a}}(\t-2)=\z_{12}\vartheta_{\bm{\rho_{\ell_1,\ell_2}},-\bm{a}}(\t).
\end{equation*}
 Applying \eqref{ref3} with $\t$ replaced by $-\frac{1}{\t}$ and then using Proposition \ref{prop1} gives the second transformation.\qedhere 
%{\bf KB: In the details file until the final displayed equation.} As inversion is an involution this gives the second transformation of \eqref{hathtrans}. Precisely,
\end{proof}

\end{document}
